\begin{table*}[t]
\centering
\caption{Posterior climate perturbation effects on transmission coefficients under the best-supported climate-driven model, with each climate covariate shifted by one standard deviation of its weekly series, so that the three covariates are compared on the same relative scale. Values show the posterior median percentage change in the time-averaged coefficient, with 95\% credible intervals, relative to the unperturbed historical-climate baseline (1000 posterior draws per scenario; excluded draws, if any, are reported per-scenario in the accompanying CSV).}
\label{tab:direct_beta_effects_posterior_1sd}
\resizebox{\textwidth}{!}{%
\begin{tabular}{llcc}
\toprule
Scenario & Direction & $\beta_R$ mean change (\%) & $\beta_H$ mean change (\%) \\
\midrule
Temperature ($\pm 1$ SD $= \pm 1.97^\circ$C) & + & -17.1 (-41.2, 13.0) & 23.3 (-7.5, 58.8) \\
Temperature ($\pm 1$ SD $= \pm 1.97^\circ$C) & - & 4.4 (-23.6, 45.3) & -18.9 (-37.0, 8.1) \\
Vapour pressure ($\pm 1$ SD $= \pm 4.32$ hPa) & + & 10.5 (-22.5, 61.8) & -56.9 (-68.3, -38.5) \\
Vapour pressure ($\pm 1$ SD $= \pm 4.32$ hPa) & - & -9.5 (-38.2, 29.0) & 131.8 (62.7, 215.3) \\
Rainfall ($\pm 1$ SD $= \pm 224$ mm) & + & -89.2 (-93.6, -84.8) & 117.3 (87.9, 156.4) \\
Rainfall ($\pm 1$ SD $= \pm 224$ mm) & - & 826.5 (557.0, 1473.5) & -54.0 (-61.0, -46.8) \\
\bottomrule
\end{tabular}%
}
\end{table*}